\section{Gruppearbejde}
Gruppen har som tidligere nævnt arbejdet ud fra en Scrum-lignende proces gennem hele projektforløbet. I denne struktur har gruppemedlemmerne haft forskellige roller, som både har været baseret på personlige styrker og på deres samarbejdsevner.
\\\\
Nikolaj har fungeret som gruppens Scrum Master gennem hele projektet. Denne rolle har passet godt til hans strukturelle tilgang til arbejde samt hans erfaring med planlægning og ledelse fra tidligere ophold i forsvaret. Han har haft fokus på at sikre overblik over arbejdsprocessen, holde styr på sprintforløb og understøtte gruppens fremdrift. Derudover har han bidraget til at skabe struktur i arbejdet og sikre, at opgaver blev fordelt og fulgt op på.
\\\\
Victor har haft rollen som Product Owner. Denne rolle er blevet tildelt ham på baggrund af hans stærke kommunikative evner og hans evne til at formidle idéer klart til både gruppen og eksterne parter. Victor har desuden stået for den primære kommunikation med projektets “kunde”, hvilket har været vigtigt for at sikre, at projektets retning og krav blev forstået og opdateret løbende gennem forløbet. Hans rolle har derfor været central i forhold til at afstemme forventninger og sikre, at projektet bevægede sig i den ønskede retning.
\\\\
Jonathan har fungeret som gruppens analytiske specialist. Han har bidraget med velovervejede analyser, refleksioner og faglige vurderinger i både udviklingen af spillet og i rapportskrivningen. Hans tilgang til arbejdet har i høj grad været præget af en analytisk og struktureret tankegang, som stemmer overens med hans blå DISC-personprofil. Dette har været en styrke i forhold til at sikre kvalitet i gruppens beslutninger og skabe et solidt grundlag for de tekniske og designmæssige valg i projektet.
\\\\
Gruppen har arbejdet i et Scrum-lignende system, hvor der blandt andet er blevet anvendt Planning Poker til estimering af use cases og tildeling af story points. Dette har hjulpet gruppen med at prioritere opgaver og skabe et bedre overblik over arbejdsbyrden i de enkelte sprint. Projektet har været opdelt i korte sprintforløb, typisk bestående af intensive arbejdsperioder fra mandag til ondsdag, og så fra ondsdag til fredag samt et længere “weekend-sprint” fra fredag til mandag. Denne struktur gjorde det muligt at arbejde fokuseret i afgrænsede perioder og løbende evaluere fremdriften.
\\\\
Hvert sprint blev afsluttet med et retrospektiv og et review-møde. Under retrospektivet blev det gennemførte arbejde gennemgået, og der blev “brændt” story points for færdiggjorte opgaver, hvilket gav gruppen et overblik over fremdriften og arbejdsindsatsen. I review-fasen blev arbejdsprocessen, samarbejdet og de anvendte metoder diskuteret, så gruppen løbende kunne justere og forbedre sin arbejdsgang. Dette har været en vigtig del af den iterative udvikling, da det har givet mulighed for kontinuerlig forbedring af både struktur og samarbejde.
\\\\
Ved sprint planning blev backloggen gennemgået, og nye use cases blev prioriteret og fordelt mellem gruppemedlemmerne. Dette sikre eden klar forventningsafstemning i begyndelsen af hvert sprint og gav en fælles forståelse af, hvilke mål der skulle opnås.
\\\\
Gruppen har dog ikke anvendt fuld Scrum-metodologi, da der eksempelvis ikke blev afholdt daglige stand-up møder. Dette var et bevidst valg, da gruppen ønskede en lettere og mere fleksibel tilgang til metoden, som bedre passede til projektets størrelse og arbejdsform. Derudover vurderede gruppen, at den løbende kommunikation internt var tilstrækkelig til at håndtere eventuelle udfordringer uden faste daglige møder. Da sprintene samtidig var relativt korte, oplevede gruppen, at der stadig var en god kontinuerlig forståelse for projektets status uden behov for daglige Scrum-møder.
\\\\
Samlet set har denne arbejdsform bidraget til en struktureret, men fleksibel proces, hvor roller, planlægning og løbende evaluering har været centrale elementer i gruppens samarbejde.